\chapter{The Template, Line by Line}\label{ch:template}

\file{template/template.cpp} is the file CPH copies for every new problem. It is strictly C++17, compiles without a
single warning under \texttt{-Wall -Wextra -Wshadow -Wpedantic}, and was tested three ways: as CPH builds it locally,
as a C++17 judge builds it, and as a C++20 judge builds it. This chapter walks through it top to bottom.
Your original template is kept as \file{template/template.cpp.bak}.

\section{Header}

\begin{code}
/**
 *    author:  Kshitish Madbhavi
 *    problem: $name$
 *    source:  $group$
 *    url:     $url$
 *    created: $date$ $time$
 */
\end{code}

The \texttt{\$...\$} words are CPH template variables. When CPH creates a file it replaces each one with the
corresponding field of the problem it received, so every solution records where it came from. Files made with
\key{Ctrl+Shift+N} are filled the same way by \file{tools/cp.py}; the \texttt{cptemplate} snippet fills them with
VS Code's date and file-name variables.

\section{Includes and the debug hook}

\begin{code}
#include <bits/stdc++.h>
using namespace std;

#ifdef LOCAL
#include "debug.h"
#else
#define dbg(...)
#endif
\end{code}

\cc{LOCAL} is defined by every local build (CPH, F5, the stress tester) and never by a judge. Locally,
\file{include/debug.h} provides \cc{dbg(...)}; on the judge \cc{dbg(...)} is defined as nothing, so
\cc{dbg(a, b);} becomes an empty statement. Chapter~\ref{ch:testing} shows what \cc{dbg} prints.

\section{Types}

\begin{code}
#define int long long

using ll = long long;
using ld = long double;
using pii = pair<int, int>;
using pll = pair<ll, ll>;
using vi = vector<int>;
using vvi = vector<vi>;
template <class T> using minpq = priority_queue<T, vector<T>, greater<T>>;
\end{code}

\cc{#define int long long} turns every \cc{int} in your code into a 64-bit integer, which removes the most common
source of wrong answers: overflow of intermediate values such as \cc{a * b} when both are up to $10^9$.

\begin{fact}[Fixed from the old template]
In the old template \cc{#define int long long} came \emph{after} the aliases, so \cc{pii} was still
\cc{pair<int,int>} with 32-bit ints. Pushing a \cc{long long} distance into a \cc{minpqpii} silently truncated it,
a classic source of mysterious Dijkstra failures. The define now comes first, so every alias is 64-bit.
\end{fact}

Consequences of \cc{#define int long long} to remember:
\begin{itemize}
  \item \cc{main} must be declared as \cc{signed main()} (it is), because \cc{long long main()} is illegal.
  \item \cc|printf("%d", x)| is wrong for an \cc{int}; use \cc{cout} or \cc|%lld|.
  \item \cc{max(x, 0)} fails to compile when \cc{x} is \cc{int}, because \cc{0} is a real 32-bit \cc{int}
        literal. Write \cc{max<int>(x, 0)}, \cc{max(x, 0LL)} or use \cc{chmax} below (it accepts mixed types).
  \item \cc{__builtin_popcount(x)} truncates to 32 bits; use \cc{__builtin_popcountll}.
  \item Memory doubles for \cc{vector<int>}. $10^7$ \cc{int}s are now 80\,MB instead of 40\,MB.
\end{itemize}

\cc{minpq<T>} is a min-heap: \cc|minpq<pii> pq; pq.push({dist, v});| pops the smallest pair first. It replaces the
four separate \cc{minpqi}/\cc{minpql}/\cc{minpqpii}/\cc{minpqpll} aliases of the old template.

\section{Macros}

\begin{code}
#define all(x) (x).begin(), (x).end()
#define rall(x) (x).rbegin(), (x).rend()
#define sz(x) (int)(x).size()
#define rep(...)   // rep(i, n): 0..n-1   rep(i, a, b): a..b-1
#define per(...)   // same ranges, backwards
\end{code}

\begin{center}
\begin{tabularx}{\linewidth}{>{\ttfamily}p{4.2cm}>{\ttfamily}p{5.8cm}L}
\toprule
\thd{You write} & \thd{It means} & \thd{Typical use} \\
\midrule
sort(all(v)) & sort(v.begin(), v.end()) & any STL algorithm on a whole container \\
sort(rall(v)) & sort(v.rbegin(), v.rend()) & sort in decreasing order \\
sz(v) & (int)v.size() & avoids signed/unsigned warnings and \texttt{v.size() - 1} underflow \\
rep(i, n) & for (int i = 0; i < n; i++) & the everyday loop \\
rep(i, 1, n + 1) & for (int i = 1; i < n + 1; i++) & 1-indexed loops \\
per(i, n) & for (int i = n - 1; i >= 0; i--) & reverse loops \\
per(i, a, b) & for (int i = b - 1; i >= a; i--) & reverse over $[a, b)$ \\
\bottomrule
\end{tabularx}
\end{center}

\cc{rep} picks its two- or three-argument form by counting arguments (the \cc{PICK3} helper macro), a standard
C++17-compatible preprocessor trick.

\begin{pitfall}[Why \texttt{sz} matters]
\cc{v.size()} is unsigned. If \cc{v} is empty, \cc{v.size() - 1} is $2^{64}-1$, so
\cc{for (int i = 0; i < v.size() - 1; i++)} runs forever. \cc{sz(v) - 1} is $-1$ and the loop is skipped.
Library code avoids the \cc{sz} name for anything else, because the macro would swallow it.
\end{pitfall}

\section{Constants}

\begin{code}
constexpr ll INF = 4e18;
constexpr int MOD = 1'000'000'007;  // 998'244'353
\end{code}

\cc{INF} is large enough to mean ``unreachable'' yet \cc{INF + INF} still fits in a \cc{long long}
($8\times10^{18} < 9.22\times10^{18}$). \cc{MOD} is used by every modular template in the library (\texttt{cpmint},
\texttt{cpncr}, \ldots); switch the comment when a problem wants 998\,244\,353.

\section{Helpers}

\begin{code}
chmin(a, b);  chmax(a, b);   // a = min(a, b) / max(a, b); true if a changed; mixed types OK
cdiv(a, b);   fdiv(a, b);    // ceil(a / b) and floor(a / b), correct for negative numbers too
cin >> v;                    // reads a whole vector (sized beforehand) or a pair
cout << v;                   // prints a vector separated by spaces
print(a, b, c);              // prints "a b c" and a newline
yn(ok);                      // prints YES or NO
mex(v);                      // smallest non-negative integer not in v, O(n)
Mex m(n); m.add(x); m.remove(x); m.get();   // mex of a changing multiset (at most n values), O(log n)
\end{code}

\begin{center}
\begin{tabular}{>{\ttfamily}l>{\ttfamily}r>{\ttfamily}r>{\ttfamily}r}
\toprule
\thd{expression} & \thd{C++ \texttt{/}} & \thd{cdiv} & \thd{fdiv} \\
\midrule
7 / 2 & 3 & 4 & 3 \\
-7 / 2 & -3 & -3 & -4 \\
7 / -2 & -3 & -3 & -4 \\
-8 / -3 & 2 & 3 & 2 \\
\bottomrule
\end{tabular}
\end{center}

C++ division truncates toward zero, which is neither floor nor ceiling for negative results; \cc{cdiv} and
\cc{fdiv} were checked on all sign combinations. \cc{chmin}/\cc{chmax} take two type parameters so
\cc{chmax(ans, 0)} compiles even though \cc{0} is a plain \cc{int}.

\section{Automatic USACO file I/O}\label{sec:fileio}

\begin{code}
void file_io() {
#ifndef LOCAL
    auto file = [](const string& spec) { /* extract "fileName":"..." */ };
    string in = file(R"($input$)"), out = file(R"($output$)");
    if (!in.empty()) freopen(in.c_str(), "r", stdin);
    if (!out.empty()) freopen(out.c_str(), "w", stdout);
#endif
}
\end{code}

\begin{figure}[H]
\centering
\begin{tikzpicture}[
  bx/.style={draw=cpblue, fill=cpblue!5, rounded corners=3pt, thick, align=left, font=\ttfamily\scriptsize, inner sep=5pt},
  arr/.style={->, very thick, cpteal}, lab/.style={font=\sffamily\scriptsize, text=cpgray, align=center}]
  \node[bx] (json) at (0, 0) {"input": \{"type": "file",\\ \ \ \ \ \ \ \ \ \ "fileName": "moocast.in"\}\\
    "output": \{"type": "file",\\ \ \ \ \ \ \ \ \ \ \ "fileName": "moocast.out"\}};
  \node[bx, draw=cpteal, fill=cpteal!5] (tmpl) at (0, -2.6) {string in = file(R"(\$input\$)");};
  \node[bx, draw=cporange, fill=cporange!6] (res) at (8.1, -2.6) {string in = file(R"("type":"file",\\
    \ \ \ \ \ \ \ \ \ \ \ \ \ \ \ \ \ \ "fileName":"moocast.in")");};
  \node[bx, draw=cpgreen, fill=cpgreen!6] (judge) at (8.1, 0) {judge: freopen("moocast.in") ...\\locally (-DLOCAL): skipped, CPH uses stdin};
  \node[lab, above=2pt of json] {Competitive Companion sends};
  \node[lab, below=2pt of tmpl] {\file{template.cpp}};
  \draw[arr] (json) -- node[left, lab] {CPH} (tmpl);
  \draw[arr] (tmpl) -- node[above, lab] {\$input\$ replaced} (res);
  \draw[arr] (res) -- node[right, lab] {at run time} (judge);
\end{tikzpicture}
\caption{How a USACO problem switches on file I/O without any manual step.}
\end{figure}

CPH replaces \texttt{\$input\$} with the problem's input description, minus its outer braces. For USACO that is
\texttt{"type":"file","fileName":"moocast.in"}, so \cc{file()} finds a file name and \cc{freopen} runs; for every
other judge it is \texttt{"type":"stdin"}, no file name is found and nothing happens. The raw string literal
\cc{R"(...)"} means the quotes need no escaping. The whole function is compiled only when \cc{LOCAL} is
\emph{not} defined, so CPH keeps running your samples through standard input.

\begin{tip}
If a file was created some other way and you need file I/O, change the placeholder by hand to
\texttt{R"("fileName":"name.in")"} (and the output one to \texttt{name.out}), or recreate the file with
\key{Ctrl+Shift+U}.
\end{tip}

\section{solve() and main()}

\begin{code}
void solve() {

}

signed main() {
    ios::sync_with_stdio(false);
    cin.tie(nullptr);
    cout << fixed << setprecision(10);
    file_io();

    int t = 1;
    // cin >> t;
    while (t--) solve();
}
\end{code}

\begin{itemize}
  \item \cc{sync_with_stdio(false)} and \cc{cin.tie(nullptr)} make \cc{cin}/\cc{cout} as fast as \cc{scanf}. After them,
        do not mix \cc{cin} with \cc{scanf} or \cc{cout} with \cc{printf}.
  \item \cc{fixed << setprecision(10)} prints floating-point answers with ten decimals, which satisfies every
        ``absolute or relative error $10^{-6}$'' checker. Integers are unaffected.
  \item Print \cc|'\n'|, not \cc{endl}: \cc{endl} flushes the output every time and can turn a fast program into
        a TLE. The one exception is interactive problems, which need the flush (\texttt{cpinteractive}).
\end{itemize}
